\section{Exact finite cubic certificates}\label{app:cubic-certificates}
Within each group, all binary vectors with the same success count have
the same polynomial value. The certificates therefore work with counts:
enumerating count vectors covers every vertex, and hence every
dependence pattern relevant to the envelope, without enumerating all
$2^{192}$ binary vectors.

For the six-term form \eqref{eq:cubic-six-family}, let
$\Phi_m(A,B,C)$ be its binary count polynomial as defined in the proof of
Proposition~\ref{prop:cubic-finite}. Table~\ref{tab:cubic-duals} gives
integers $D,b_A,b_B,b_C,b_0$ such that
\begin{equation}\label{eq:cubic-integer-dual}
 D\Phi_m(A,B,C)\ge b_AA+b_BB+b_CC+b_0
       \quad(0\le A,B,C\le m\text{ integers}).
\end{equation}
Table~\ref{tab:cubic-primals} gives attaining laws. Each listed weight is
nonnegative, the weights sum to one, and their mean count vector is
$(m/4,m/2,3m/4)$. Choosing uniform success subsets conditional on these
counts therefore gives the required singleton marginals. The checker
below verifies, in rational arithmetic, the inequality at every integer
count triple and the equality of primal and dual expectations.

\begin{table}[htbp]
\centering
\begin{tabular}{rrrrrr}
\toprule
$m$&$D$&$b_A$&$b_B$&$b_C$&$b_0$\\\midrule
6&13&962&778&842&$-4816$\\
8&7&793&770&798&$-5882$\\
64&105&871710&900446&899046&$-51743768$\\\bottomrule
\end{tabular}
\caption{Integer affine minorants for the three-group witnesses.}
\label{tab:cubic-duals}
\end{table}

\begin{table}[htbp]
\centering
\begin{tabular}{rll}
\toprule
$m$&Count states $(A,B,C)$&Corresponding probabilities\\\midrule
6&$(1,4,4),(2,1,6),(6,2,1)$&$17/26,\ 4/13,\ 1/26$\\
8&$(1,2,8),(1,5,6),(8,4,2)$&$2/7,\ 4/7,\ 1/7$\\
64&$(4,19,64),(5,19,64),(16,40,43),(64,31,17)$
 &$241/735,\ 12/735,\ 419/735,\ 63/735$\\\bottomrule
\end{tabular}
\caption{Primal laws attaining the affine bounds.}
\label{tab:cubic-primals}
\end{table}

For $m=8$, minimizing the residual in
\eqref{eq:cubic-integer-dual} over $A,B$ for each $C=0,\ldots,8$ gives
\[
 168,\ 42,\ 0,\ 30,\ 22,\ 0,\ 0,\ 16,\ 0.
\]
The three primal states have values $405,507,734$, respectively, with
expected value $3572/7$. In all three examples the orbit sizes and the
upper and lower values per unweighted monomial are
\begin{align*}
 s&=\left(\binom m3,m\binom m2,\binom m2,m^2,m^2,\binom m2\right),\\
 u&=(3/4,1/2,1/2,1/4,1/4,1/4),\qquad
 l=(1/4,0,0,0,0,0).
\end{align*}
Thus $\cav f=\sum_ic_is_iu_i$ and $\sum_e\vex f_e=\sum_ic_is_il_i$.
These formulas prove the remaining columns of
Table~\ref{tab:cubic-finite}. The exact hull gaps for dimensions
18, 24, 192 are $10411/52$, $3225/7$, $27562048/105$, respectively.

For the two-level $m=16$ certificate
\eqref{eq:two-level-16-dual}, the integer residual is
\[
 2\left[A\binom C2+20\binom A2\right]-428A-177C+3372.
\]
Its minimum over $C=0,\ldots,16$ for successive $A=0,\ldots,16$ is
\[
\begin{gathered}
 540,\ 352,\ 204,\ 96,\ 28,\ 0,\ 9,\ 7,\ 0,\\
 0,\ 19,\ 63,\ 132,\ 235,\ 369,\ 540,\ 742.
\end{gathered}
\]
At counts $(8,12)$ the value and its affine bound both equal 1088.

For the smaller members of \eqref{eq:two-level-family}, write
$\Phi_m(A,C)=A\binom C2+(5m/4)\binom A2$.
Table~\ref{tab:two-level-small} gives affine minorants valid at every
integer pair $0\le A,C\le m$. At the listed counts each minorant equals
$\Phi_m$. Choosing uniform success subsets of those fixed sizes,
independently in the two groups, gives singleton means $1/2$ and $3/4$
and attains the bound. Taking expectations of the minorant under any
feasible law gives the matching lower bound, which proves the stated
exact convex-envelope value. Every term has
lower envelope zero and upper envelope $1/2$, so
$\cav f_m=T(f_m)=(9m/8)\binom m2$; subtraction gives the exact ratios.
The checker verifies all $5^2+9^2+13^2$ minorant inequalities and these
attainment and envelope equalities in rational arithmetic.

\begin{table}[htbp]
\centering
\begin{tabular}{rllrrl}
\toprule
$m$&Affine minorant&Counts $(A,C)$&$\vex f_m$&$\cav f_m=T$&$T/H$\\\midrule
4&$9A+4C-19$&$(2,3)$&11&27&$27/16$\\
8&$47A+20C-188$&$(4,6)$&120&252&$21/11$\\
12&$(231/2)A+48C-684$&$(6,9)$&441&891&$99/50$\\\bottomrule
\end{tabular}
\caption{Exact smaller two-level certificates and attaining counts.}
\label{tab:two-level-small}
\end{table}

\newpage
The following complete Python 3 checker uses only the standard library.
It enumerates $7^3+9^3+65^3$ three-group states and
$5^2+9^2+13^2+17^2$ two-group states, using exactly the data in the
tables; no numerical optimization or solver is involved. Concatenating
the two code blocks in order gives the supplied file
\texttt{verification/check\_stage02\_finite.py}.
\par\medskip\noindent\begin{minipage}{\textwidth}\small
\begin{verbatim}
from fractions import Fraction as Q
from itertools import product
from math import comb

# m, coefficients, denominator, integer affine coefficients,
# probability denominator, (state, probability numerator), ratio
cases = [
 (6, (2,3,9,10,7,7), 13, (962,778,842,-4816), 26,
  [((1,4,4),17), ((2,1,6),8), ((6,2,1),1)],
  Q(20891,10411)),
 (8, (2,3,13,12,8,7), 7, (793,770,798,-5882), 7,
  [((1,2,8),2), ((1,5,6),4), ((8,4,2),1)],
  Q(6601,3225)),
 (64, (2,3,120,105,70,63), 105,
  (871710,900446,899046,-51743768), 735,
  [((4,19,64),241), ((5,19,64),12),
   ((16,40,43),419), ((64,31,17),63)],
  Q(7443345,3445256))]

def phi(a, b, c, coef):
    values = (comb(c,3), b*comb(c,2), comb(b,2),
              a*c, a*b, comb(a,2))
    return sum(v*w for v,w in zip(values,coef))

\end{verbatim}
\end{minipage}
\newpage
\noindent\begin{minipage}{\textwidth}\small
\begin{verbatim}
for m,coef,D,dual,P,atoms,ratio in cases:
    for a,b,c in product(range(m+1), repeat=3):
        rhs = sum(v*w for v,w in zip((a,b,c,1),dual))
        assert D*phi(a,b,c,coef) >= rhs
    assert all(p >= 0 for _,p in atoms)
    assert sum(p for _,p in atoms) == P
    means = [sum(Q(p,P)*state[j] for state,p in atoms)
             for j in range(3)]
    assert means == [Q(m,4), Q(m,2), Q(3*m,4)]
    primal = sum(Q(p,P)*phi(*state,coef) for state,p in atoms)
    lower = sum(Q(v,D)*w for v,w in zip(dual,means+[1]))
    assert primal == lower
    sizes = (comb(m,3), m*comb(m,2), comb(m,2),
             m*m, m*m, comb(m,2))
    upper = (Q(3,4), Q(1,2), Q(1,2), Q(1,4), Q(1,4), Q(1,4))
    cav = sum(c*s*u for c,s,u in zip(coef,sizes,upper))
    term_lower = Q(coef[0]*comb(m,3),4)
    assert cav > primal
    assert (cav-term_lower)/(cav-primal) == ratio
    assert ratio > 2
    print(3*m, cav, term_lower, primal, ratio)

# m, affine coefficients (A,C,1), attaining value, cav=T, ratio
small = [
 (4, (9,4,-19), 11, 27, Q(27,16)),
 (8, (47,20,-188), 120, 252, Q(21,11)),
 (12, (Q(231,2),48,-684), 441, 891, Q(99,50))]
for m,dual,vex,cav,ratio in small:
    def value(a,c):
        return a*comb(c,2)+Q(5*m,4)*comb(a,2)
    slacks = [value(a,c)-dual[0]*a-dual[1]*c-dual[2]
              for a,c in product(range(m+1), repeat=2)]
    assert min(slacks) == 0
    a,c = m//2,3*m//4
    assert Q(a,m) == Q(1,2) and Q(c,m) == Q(3,4)
    assert value(a,c) == dual[0]*a+dual[1]*c+dual[2] == vex
    assert (m+Q(5*m,4))*comb(m,2)/2 == cav
    assert cav > vex and Q(cav,cav-vex) == ratio < 2
    print(2*m, cav, vex, ratio)

mins = []
for a in range(17):
    residuals = [2*(a*comb(c,2)+20*comb(a,2))
                 -428*a-177*c+3372 for c in range(17)]
    mins.append(min(residuals))
assert mins == [540,352,204,96,28,0,9,7,0,0,19,63,
                132,235,369,540,742]
assert 8*comb(12,2)+20*comb(8,2) == 1088
assert 214*8+Q(177,2)*12-1686 == 1088
assert Q(2160,2160-1088) == Q(135,67)
assert Q(135,67) > 2
assert Q(243,115)*(1-Q(1,20)) == Q(4617,2300) > 2
print("All finite cubic certificates passed.")
\end{verbatim}
\end{minipage}
